% called by main.tex
%
\chapter{Introduction}
\label{ch::chapter1}
En los últimos años, los dispositivos IoT han ganado una importancia significativa tanto en la industria como en la vida cotidiana, impulsados por el crecimiento del número de usuarios de Internet. Se estima que la cantidad de dispositivos IoT alcanzará los 18 mil millones en los próximos años, con proyecciones que indican que para 2030 podrían superar los 40 mil millones \cite{kumar_how_2024}.

Paralelamente, la computación cuántica ha experimentado avances significativos. Estas nuevas arquitecturas computacionales, que aprovechan los principios de la mecánica cuántica, son capaces de resolver problemas matemáticos de alta complejidad en tiempos inalcanzables para los ordenadores clásicos. Como consecuencia, los sistemas de cifrado ampliamente utilizados en la actualidad, como RSA y ECC, podrían volverse inseguros ante la capacidad de los futuros ordenadores cuánticos de resolver en poco tiempo problemas que antes se consideraban intratables.

Ante esta amenaza, el \textbf{Instituto Nacional de Estándares y Tecnología (NIST, por sus siglas en inglés)} inició un proceso de estandarización de algoritmos criptográficos resistentes a ataques cuánticos, con el objetivo de garantizar la seguridad de la información en el futuro \cite{computer_security_division_post-quantum_2017}.

La seguridad en los dispositivos IoT es un desafío crítico, ya que muchos de estos dispositivos dependen de algoritmos criptográficos tradicionales que podrían volverse obsoletos. Esto representa un riesgo significativo, especialmente en entornos donde los dispositivos embebidos operan con recursos computacionales limitados y no siempre pueden ser actualizados fácilmente. En este contexto, la \textbf{criptografía post-cuántica (PQC, Post-Quantum Cryptography)} se presenta como una solución necesaria para garantizar la seguridad a largo plazo, proporcionando resistencia a ataques cuánticos sin requerir hardware extremadamente potente.

Dentro de la criptografía post-cuántica, los algoritmos pueden clasificarse en varias categorías según el problema matemático en el que basan su seguridad \cite{fernandez-carames_pre-quantum_2020}: Criptografía basada en códigos, Criptografía basada en redes euclidianas, Criptografía basada en funciones hash, Criptografía basada en isogenias de curvas elípticas, Criptografía basada en polinomios multivariantes.

De estos enfoques, los algoritmos basados en códigos han demostrado ser una alternativa robusta contra ataques cuánticos. En la última ronda del proceso de estandarización del NIST, tres algoritmos de este tipo fueron seleccionados: \textbf{HQC (Hamming Quasi-Cyclic)} \cite{noauthor_hamming_nodate}, \textbf{Classic McEliece} \cite{noauthor_classic_nodate} y \textbf{BIKE} \cite{aragon_bike_nodate}.

Este trabajo se centrará en el algoritmo \textbf{HQC}, un esquema de encapsulación de claves basado en códigos diseñado para proporcionar seguridad frente a ataques tanto de ordenadores clásicos como cuánticos \cite{noauthor_hqc_nodate-1}. En comparación con los otros esquemas finalistas, \textbf{HQC} ofrece un equilibrio entre tamaño de clave, eficiencia en el cifrado/descifrado y requisitos computacionales, lo que lo hace una alternativa interesante para dispositivos IoT. Según \cite{kuznetsov_trade-offs_nodate}, \textbf{HQC} presenta un rendimiento adecuado para dispositivos con bajas prestaciones y consumo reducido, características esenciales en muchas aplicaciones IoT.